% ==============================================================================
% CONFIGURAZIONE DEI PACCHETTI (PACKAGES) - ALGEBRA LINEARE
% ==============================================================================

\usepackage[utf8]{inputenc}
\usepackage[T1]{fontenc}
\usepackage[italian]{babel}
\usepackage{lmodern}
\usepackage{microtype}

% --- Geometria e layout della pagina ---
\usepackage[
    a4paper,
    top=2.5cm,
    bottom=2.5cm,
    left=2.7cm,
    right=2.7cm,
    headheight=14pt,
    footskip=1.2cm
]{geometry}

% --- Pacchetti matematici avanzati ---
\usepackage{amsmath}
\usepackage{amssymb}
\usepackage{amsthm}
\usepackage{mathtools}
\usepackage{amsfonts}
\usepackage{mathrsfs}
\usepackage{bm}
\usepackage{cancel}
\usepackage{siunitx}
\usepackage{nicematrix} % Per matrici avanzate, blocchi e linee tratteggiate
\usepackage{systeme}    % Per la formattazione pulita dei sistemi lineari

% --- Gestione grafica, TikZ e diagrammi commutativi ---
\usepackage{graphicx}
\usepackage{xcolor}
\usepackage{tikz}
\usepackage{tikz-cd}    % Diagrammi commutativi tra spazi vettoriali
\usetikzlibrary{calc,arrows.meta,positioning,patterns,decorations.pathmorphing,backgrounds,shapes.geometric,matrix,3d}
\usepackage{pgfplots}
\pgfplotsset{compat=1.18}

% --- Strutture, tabelle e liste ---
\usepackage{float}
\usepackage{booktabs}
\usepackage{tabularx}
\usepackage{array}
\usepackage{multirow}
\usepackage{enumitem}
\setlist{itemsep=0.2em, topsep=0.3em}

% --- Box decorati con tcolorbox ---
\usepackage[most]{tcolorbox}

% --- Intestazioni e piè di pagina ---
\usepackage{fancyhdr}
\pagestyle{fancy}
\fancyhf{}
\fancyhead[LE,RO]{\thepage}
\fancyhead[RE]{\nouppercase{\leftmark}}
\fancyhead[LO]{\nouppercase{\rightmark}}
\renewcommand{\headrulewidth}{0.4pt}
\renewcommand{\footrulewidth}{0pt}

% Stile plain per le prime pagine di capitolo
\fancypagestyle{plain}{
    \fancyhf{}
    \fancyfoot[C]{\thepage}
    \renewcommand{\headrulewidth}{0pt}
}

% --- Definizione Palette Colori Tematici ---
\definecolor{navyblue}{RGB}{10, 45, 90}
\definecolor{coolblue}{RGB}{28, 110, 164}
\definecolor{tealaccent}{RGB}{0, 128, 128}
\definecolor{forestgreen}{RGB}{20, 115, 60}
\definecolor{warmamber}{RGB}{195, 110, 10}
\definecolor{crimsonred}{RGB}{175, 30, 45}
\definecolor{purpleaccent}{RGB}{110, 40, 140}
\definecolor{lightgray}{RGB}{210, 215, 220}
\definecolor{boxgray}{RGB}{248, 249, 250}
\definecolor{subtlegray}{RGB}{240, 242, 245}

% --- Formattazione Titoli ---
\usepackage{titlesec}
\titleformat{\chapter}[display]
  {\normalfont\bfseries\color{navyblue}}
  {\filleft\fontsize{44}{44}\selectfont\color{lightgray}CAPITOLO \thechapter}
  {1ex}
  {\titlerule\vspace{1ex}\LARGE}
  [\vspace{1ex}\titlerule]

% --- Collegamenti ipertestuali e riferimenti intelligenti ---
\usepackage[
    colorlinks=true,
    linkcolor=navyblue,
    citecolor=forestgreen,
    urlcolor=coolblue,
    bookmarksnumbered=true,
    pdfstartview=FitH
]{hyperref}
\usepackage{bookmark}
\usepackage[italian,capitalize,noabbrev]{cleveref}
